% ============================================================
% 第 25 章 直流电机的换向（教材第六篇，印刷页 270--277）
% 节名：第一节 直流电机换向的电磁现象；第二节 改善换向的方法；第三节 补偿绕组
% ============================================================
\chapter{直流电机的换向}

换向是直流电机的"专利问题"：电枢元件从一条支路进入另一条支路时，元件中的电流必须从 $+i_a$ 变为 $-i_a$，这个过程在电刷下几十微秒内完成，极易产生火花。换向不良是直流电机最主要的故障源，也是限制其容量与转速提高的根本原因。

\section{直流电机换向的电磁现象}

\subsection{换向过程}
设电刷宽度等于一个换向片宽度，则被电刷短路的那个元件从开始被短路到脱离短路所经历的时间称为\textbf{换向周期} $T_k$：
\begin{equation}
T_k=\frac{b_k}{v_k},
\end{equation}
其中 $b_k$ 为换向片宽度，$v_k$ 为换向器表面线速度。$T_k$ 通常只有几毫秒到几十微秒，而元件电流却要从 $+i_a$ 变到 $-i_a$。

\subsection{换向元件中的电动势}
换向元件本身是一个电感线圈，电流的急剧变化会产生\textbf{自感电动势}（电抗电动势）：
\begin{equation}
e_L=-L_k\frac{\mathrm{d}i}{\mathrm{d}t}\approx L_k\frac{2i_a}{T_k},
\end{equation}
其方向\textbf{阻碍电流变化}，即力图维持原电流方向——这使换向电流"延迟"改变，在电刷后刷边引起火花。

此外，换向元件位于几何中性线附近，虽然主磁场在此处为零，但\textbf{交轴电枢磁场}仍然存在，换向元件切割该磁场产生\textbf{旋转电动势}（电枢反应电动势）$e_a$，方向同样阻碍换向。

\subsection{三种换向}
\begin{itemize}
  \item \textbf{直线换向}：$e_L+e_a=0$（理想情况），换向电流随时间线性变化，无火花，最理想；
  \item \textbf{延迟换向}：$e_L+e_a>0$ 占主导，电流变化滞后，后刷边电流密度大，\textbf{易产生火花}（实际电机多为此种）；
  \item \textbf{超前换向}：人为加强换向磁场使其过量，电流变化超前，前刷边电流密度大，也会产生火花。
\end{itemize}

\subsection{火花的危害}
火花会灼伤换向器与电刷表面，加剧磨损；严重时形成\textbf{环火}（电弧跨越全部换向片），烧毁电机。因此必须把火花控制在允许范围内（国家标准按火花等级分为 1、1¼、1½、2、3 级）。

\section{改善换向的方法}

改善换向的基本思路是：\textbf{用外加磁场（换向极）产生一个与 $e_L$ 方向相反的电动势，把延迟换向拉回直线换向}。

\subsection{装设换向极}
在相邻主磁极之间的几何中性线上装设\textbf{换向极}，其绕组与电枢绕组\textbf{串联}，使换向极磁动势随电枢电流成正比变化。换向极的极性应使它在换向元件中产生的电动势 $e_k$ 与 $e_L$ 相反，即
\begin{equation}
e_k\approx-\left(e_L+e_a\right).
\end{equation}
换向极极性的确定：对\textbf{发电机}，换向极极性与电枢旋转方向上前一个主磁极的极性\textbf{相同}；对\textbf{电动机}，则\textbf{相反}。换向极磁路应不饱和，以保证 $e_k\propto I_a$ 在整个负载范围内都能补偿。

\subsection{移刷（移动电刷位置）}
把电刷从几何中性线移开一个角度，使换向元件处于主磁场中，利用主磁场产生的电动势抵消 $e_L$。发电机应\textbf{顺转向}移刷，电动机应\textbf{逆转向}移刷。缺点是只对某一负载有效，负载变化时补偿效果变差，且移刷会引入直轴电枢反应，故工程上只作为辅助手段。

\subsection{减小电抗电动势}
\begin{itemize}
  \item 减小换向元件的匝数（每元件匝数取 1）；
  \item 减小元件的电感（缩短端部、增大槽口）；
  \item 适当增大换向片数、减小片间电压。
\end{itemize}

\subsection{选用合适的电刷}
电刷的接触特性对换向影响很大：石墨电刷接触电阻大、抑制火花能力强，但允许的电流密度小；金属石墨电刷导电好、电流密度大，但换向性能差。应按电流密度与换向要求选择。

\section{补偿绕组}

\subsection{作用}
换向极只能补偿几何中性线附近的换向，无法消除交轴电枢反应对主磁场的\textbf{畸变}。当电机负载重、过载大、或采用晶闸管供电（电流脉动大）时，畸变会使极靴下一侧磁密过高、另一侧过低，导致片间电压过高而环火。为此需装设\textbf{补偿绕组}。

\subsection{结构与连接}
补偿绕组嵌放在\textbf{主磁极极靴的槽内}，其导体中电流方向与该极下电枢导体电流方向\textbf{相反}，从而直接抵消交轴电枢反应磁动势。补偿绕组与电枢绕组\textbf{串联}，以保证补偿作用随负载自动成正比变化。

\subsection{特点}
\begin{itemize}
  \item 补偿效果好，可使气隙磁场基本不畸变，允许电机做大容量、高过载；
  \item 结构复杂、造价高（极靴需开槽、嵌线），故只用于大型或特殊要求的直流电机；
  \item 补偿绕组使换向极的负担减轻，二者常配合使用。
\end{itemize}

\begin{center}
\begin{tabular}{@{}p{2.2cm}p{3.8cm}p{4.2cm}p{3.4cm}@{}}
\toprule
措施 & 针对的问题 & 原理 & 适用范围 \\
\midrule
换向极 & 换向元件中的电抗电动势 $e_L$ & 串联绕组产生反向电动势 & 几乎所有直流电机 \\
移刷 & $e_L$ 引起的延迟换向 & 利用主磁场产生补偿电动势 & 辅助手段，负载固定时 \\
减小元件电感 & $e_L$ 本身 & 减少匝数、缩短端部 & 设计阶段 \\
补偿绕组 & 交轴电枢反应造成的磁场畸变 & 极靴槽内绕组直接抵消 & 大型、重载、晶闸管供电 \\
\bottomrule
\end{tabular}
\end{center}

\corepoint{本章的因果链要背熟：\textbf{换向 = 电流急剧反向 $\rightarrow$ 电感产生 $e_L$（阻碍变化）$\rightarrow$ 换向延迟 $\rightarrow$ 后刷边火花}。改善手段都是"制造一个反向电动势"：换向极（串联、随负载自动补偿）是主力，移刷是辅助，补偿绕组解决的是另一个问题——\textbf{交轴电枢反应引起的磁场畸变}（不是换向本身）。考试常问两者的区别，务必分清"换向极管换向、补偿绕组管畸变"。}
